% P21.3 residual-Picard curve regularity, 11 September 2026.
\begin{proposition}[Residual Picard criterion for ordinary smooth harmonic germs]
\label{prop:p213-residual-picard-smooth-curve}
Let \(K\) be algebraically closed, complete, and nontrivially valued, let
\(C\) be a smooth proper algebraic \(K\)-curve, and let
\(W\subset C^{\mathrm{an}}\) be open.  Suppose that
\(h\in C^0(W,\mathbf R)\) satisfies
\[
                         \dprimeddprime[h]=0.
\]
For a type~\(2\) point \(x\in W\), let \(C_x\) be its smooth projective
residue curve and let
\[
 \operatorname{ob}_x(h)\in
 \operatorname{Pic}^0(C_x)\otimes_{\mathbf Z}\mathbf R
\]
be the residual class of the harmonic germ of \(h\) at \(x\).  Then the
following are equivalent:
\begin{enumerate}
\item \(h\) is ordinary CLD-smooth near \(x\) and
      \(\dprimeddprime h=0\) there as a smooth form;
\item \(\operatorname{ob}_x(h)=0\).
\end{enumerate}
At points of types \(1\), \(3\), and \(4\), the germ of \(h\) is
automatically ordinary CLD-smooth with zero smooth Hessian.

Consequently, if
\begin{equation}
\label{eq:p213-residual-picard-vanishing}
       \operatorname{Pic}^0(C_x)\otimes_{\mathbf Z}\mathbf R=0
       \qquad\text{for every type~\(2\) point }x\in W,
\end{equation}
then \(h\) is ordinary CLD-smooth on \(W\) and
\(\dprimeddprime h=0\) as a smooth form.  The hypothesis
\eqref{eq:p213-residual-picard-vanishing} holds, in particular, when the
residue field \(\widetilde K\) is algebraic over a finite field, or when
\(C\) is a Mumford curve.
\end{proposition}
